	\ifdefined\standalone
\documentclass[12pt]{article}
\usepackage{graphicx}
\usepackage{float}
\usepackage[a4paper, margin=1in]{geometry}
\usepackage[utf8]{inputenc}
\usepackage{textcomp}
\usepackage{amsmath}
\usepackage{booktabs}
\usepackage{tikz}
\usepackage{enumitem}
\usetikzlibrary{decorations.pathmorphing,patterns}
\usepackage[hidelinks]{hyperref}
\begin{document}
\fi

\section{Reaction Enthalpy}

This case is designed to verify the implementation of reaction enthalpy in the code. The volumetric chemical heat source term, denoted $\dot{Q}^{'''}_{sc}$ is expressed as
\begin{equation}
\dot{Q}^{'''}_{sc} = - \sum_{k}^{K} \dot{\omega}^{'''}_{r,k} \, \Delta H_{r,k}
+ \sum_{i}^{M} \dot{\omega}^{'''}_{net,i} \, h_i 
\label{eq:chimical source term}
\end{equation}
where $K$ is the number of reactions, $\dot{\omega}^{'''}_{r,k}$ is the reaction rate of reaction $k$, $\Delta H_{r,k}$ is the reaction enthalpy, $\dot{\omega}^{'''}_{net,i}$ is the net production rate of condensed species $i$, and $h_i$ is the specific sensible enthalpy of species $i$.

A one-dimensional simulation is performed, consisting of a single condensed material, denoted material A. This material undergoes a single reaction transforming it into material B without producing gaseous species. The reaction is first order, with an activation energy of $E = 150$~kJ\,mol$^{-1}$ and a pre-exponential factor of $5 \times 10^{13}$~s$^{-1}$. Two sub-cases are considered. The first corresponds to an endothermic reaction with $\Delta H_k = 1 \times 10^{6}$~J\,kg$^{-1}$, while the second corresponds to an exothermic reaction with $\Delta H_k = -1 \times 10^{5}$~J\,kg$^{-1}$. A third case is also investigated, in which the reaction is slightly modified to produce gaseous species. In this configuration, $60\%$ of the mass of material A is converted into gases, and the reaction enthalpy is set to $\Delta H_k = 0$. As a result, the first term in Eq.~\ref{eq:chimical source term} vanishes, and only the second term, corresponding to the variation of sensible enthalpy of the condensed phase, contributes to the heat source term.

The numerical Chemical source term and their comparison with analytical solutions are presented in Figure~\ref{fig:endo_results} for the endothermic case, Figure~\ref{fig:exo_results} for the exothermic case, and Figure~\ref{fig:sensible_results} for the sensible enthalpy variation case. The analytical solutions are computed directly from Eq.~\ref{eq:chimical source term}, using reaction rates extracted from the numerical solution.


\begin{figure}[H]
\centering
\begin{minipage}[t]{0.49\textwidth}
\centering
	\includegraphics[width=\linewidth]{endothermic_reaction/Reaction_Enthalpy_Endo.png}
    \caption{Endothermic reaction case.}
\label{fig:endo_results}
\end{minipage}
\hfill
\begin{minipage}[t]{0.49\textwidth}
\centering
    \includegraphics[width=\linewidth]{exothermic_reaction/Reaction_Enthalpy_Exo.png}
    \caption{Exothermic reaction case.}
\label{fig:exo_results}
\end{minipage}
\hfill
\begin{minipage}[t]{0.49\textwidth}
\centering
    \includegraphics[width=\linewidth]{sensible_enthalpy/Reaction_Enthalpy_sensible.png}
    \caption{Sensible enthalpy case.}
\label{fig:sensible_results}
\end{minipage}
\end{figure}


\ifdefined\standalone
\end{document}
\fi




